\section{Conclusion}
\label{sec:conclusion}

In this paper, we presented \textbf{RelieF-CR} (Relief-Guided Cloud Removal), a comprehensive multi-modal cross-attention generative framework for high-resolution ($5.0\,\text{m}$) satellite optical cloud removal. By synergistically integrating four heterogeneous sensor modalities---ISRO LISS-IV optical imagery, Sentinel-1 dual-polarization SAR, Sentinel-2 dry-season temporal references, and CartoDEM 4-channel topographic surface models---the proposed architecture effectively resolves the physical underdetermination of SAR-to-optical translation.

Key technical innovations include:
\begin{enumerate}
    \item Learned Spatial Transformer Networks (STNs) that rectify inter-sensor registration jitter.
    \item Windowed Multi-Head Cross-Attention that adaptively weights and integrates radar and temporal structural representations in degraded optical regions.
    \item A dual-head heteroscedastic generator providing calibrated, pixel-wise uncertainty maps ($\rho = +0.245$).
    \item A 7-term loss suite incorporating frequency-domain FFT and NDVI vegetation consistency constraints.
\end{enumerate}

Comprehensive experimental evaluation across held-out test patches demonstrates state-of-the-art performance with \textbf{30.96\,dB PSNR}, \textbf{0.965 SSIM}, \textbf{1.55$^\circ$ SAM}, and \textbf{7.29 ERGAS}. Exhaustive diagnostic auditing mathematically confirms that our model's structural correlation with SAR ($0.4791$) matches real optical ground truth ($0.4766$), ruling out radar over-imprinting artifacts. Finally, memory-mapped streaming inference enables continuous, seamless reconstruction of massive $300\,\text{MPix}$ full-swath satellite scenes, paving the way for reliable, year-round automated Earth observation analytics.
